\documentclass[12pt]{article}
\usepackage{amsmath, amssymb, amsthm}
\usepackage{parskip}
\usepackage{setspace}
\usepackage{titlesec}

\onehalfspacing
\setlength{\jot}{12pt}
\titleformat{\section}{\large\bfseries}{\thesection}{1em}{}
\titleformat{\subsection}{\normalsize\bfseries}{\thesubsection}{1em}{}

\theoremstyle{definition}
\newtheorem{definition}{Definition}

% For the space in definitions
\makeatletter
\def\thm@space@setup{%
  \thm@preskip=1.5em
  \thm@postskip=1.5em
}
\makeatother

\theoremstyle{remark}
\newtheorem*{remark}{Remark}

\title{Why Measure Theory}
\author{Carter Deslauriers}
\date{\today}

\begin{document}
\maketitle

\section{Introduction: Who Cares}

Write in the Introduction about why measure in general


\section{Pointwise Convergence: Breaking at the Limit}
\subsection{Motivation for DCT}
\subsubsection{Pointwise Convergence}

% EXAMPLE FUNCTION
Consider the sequence of functions defined on $[0, 1]$

\begin{equation}
f_n(x) = \begin{cases}
    n: & 0 < x \leq \frac{1}{n} \\
    0: & x = 0, \text{ } \frac{1}{n} < x \leq 1
          \end{cases}
\end{equation}

First to show that $f_n(x) \to f(x) = 0$ pointwise consider the definition.



% DEFINITION OF POINTWISE CONVERGENCE
\begin{definition}[Pointwise Convergence]
For every $\epsilon > 0$ there exists some $N \in \mathbb{N}$ such that whenever $n \geq N$ it
follows that
\begin{equation*}
    |f_n(x) - f(x)| < \epsilon
\end{equation*}
\end{definition}

% PROOF OF POINTWISE CONVERGENCE
To prove the function converges pointwise I will show it satisfies the defintion

\begin{proof}
Subsituting our functions gives

\begin{equation*}
    |f_n(x) - 0| = 0 < \epsilon \text{ whenever } \frac{1}{n} < x
\end{equation*}

By the archimedean property we know for any $x \in \mathbb{R}$ there exists an $N \in \mathbb{N}$
such that whenever

\begin{align*}
    n \geq N
    &\implies \frac{1}{n} < x \\
    &\implies |f_n(x) - 0| = 0 \\
    &\implies |f_n(x) - 0| < \epsilon
\end{align*}

For the special case of $x = 0$ observe that $f_n(0) = 0$ giving 

\begin{align*}
    f_n(0) = 0
    &\implies |f_n(0) - 0| = 0 \\
    &\implies |f_n(0) - 0| < \epsilon
\end{align*}

Meaning our function converges pointwise to $f(x)$ for $f(x) = 0$

\end{proof}


Now to show the function does not converge uniformly first consider the
defintion of uniform convergence

% DEFINITION OF UNIFORM CONVERGENCE
\begin{definition}[Uniform Convergence]
    For every $\epsilon > 0$ there exists some $N \in \mathbb{N}$ such that whenever $n \geq N$ it
    follows that for all $x$
    \begin{equation*}
        |f_n(x) - f(x)| < \epsilon
    \end{equation*}
\end{definition}

To show that $f_n(x) \not\to f(x)$ uniformally, I will show there exists no such $N$
satisfying the defintion for all $x$ at once.


% PROOF FOR NON UNIFROM CONVERGENCE
\begin{proof}

For contradiction assume there exists $N \in \mathbb{N}$ such that whenever $n \geq N$
\begin{equation*}
    |f_n(x) - f(x)| < \epsilon
\end{equation*}
for all $x$

Observe that for any such $N$

\begin{equation*}
    f_n(x) = n, \text{ for } 0 < x \leq\frac{1}{n}
\end{equation*}

Choose any $\epsilon \leq 1$ specifcally I choose $\frac{1}{2}$ and call it $\epsilon_0$, 
then it is that case that for $0 < x \leq\frac{1}{n}$

\begin{equation*}
    |f_n(x) - 0| = n \implies |f_n(x) - 0| \geq \epsilon_0
\end{equation*}

But thats a contradiction since $f_n(x) \to f(x)$ uniformally gives an $N$ that works
for all $x$ at once. Meaning $f_n(x) \not\to f(x)$ uniformally.

\end{proof}


% BREAKING LIMIT SECTION
\subsubsection{Limit Breaking}

To show that

\begin{equation}
    \lim_{n \to \infty}\int_0^1 f_n(x)\,dx \neq \int_0^1 f(x)\,dx
\end{equation}

Consider both methods of integration and there definitions

% RIEMANN INTEGRATION DEFINTION
\begin{definition}[Reimann Integration]
    Let $f$ be a bounded function $f: [a, b] \to \mathbb{R}$.
    Let $P$ be a partition of the domain of the form
    \begin{equation*}
        P = \{x_0, x_1, \ldots, x_n\}
    \end{equation*}
    Such that 
    \begin{equation*}
        a = x_0 < x_1 < \ldots < x_n = b
    \end{equation*}
    If for every $\epsilon > 0$ there exists a partion $P_e$ 
    such that $U(f, P_\epsilon) - L(f, P_\epsilon) < \epsilon$ then $f(x)$ is
    Riemann Integrable with $\int f = U(f) = L(f).$

    Where 
    \begin{equation*}
        U(f, P_\epsilon) = \sum_{k=1}^{n} (x_k - x_{k-1}) \sup_{[x_{k-1}, x_k]} f
    \end{equation*}
    \begin{equation*}
        L(f, P_\epsilon) = \sum_{k=1}^{n} (x_k - x_{k-1}) \inf_{[x_{k-1}, x_k]} f
    \end{equation*}
    
    and
    \begin{equation*}
        U(f) = \inf_P U(f, P)
    \end{equation*}
    \begin{equation*}
        L(f) = \sup_P L(f, P)
    \end{equation*}

\end{definition}


% LEBESEGUE INTEGRATION DEFINITION
\begin{definition}[Lebesgue Integration]
    Let $(X, S, \mu)$ be a measure space and let $f(x) \to \mathbb{R}$
    then $\int f = \sup L(f, P)$.

    Where P is an $S$-Partition. Where an $S$-Partition is of the form
    \begin{equation*}
        P = {E_1, E_2, \ldots, E_n}
    \end{equation*}
    Where each $E_k \in S$ and 
    \begin{equation*}
        \bigcup_k E_k = X
    \end{equation*}
    
    and
    \begin{equation*}
        \mathbb{L}(f, P) = \sum_k \mu(E_k)\inf_{E_k} f
    \end{equation*}


\end{definition}

Now I will show in accordance with the definitions neither integration method
will suffice swapping the integral and the limit

% PROOF RIEMANN INTEGRATION FAILS AT THE LIMIT
\begin{proof}
First I will show that for any fixed n $f_n(x)$ is Riemann Integrable.

Fix $n \in \mathbb{N}$, first observe $|f_n(x) \leq n$
for all $x \in [0, 1]$. Now take the partition where
\begin{equation*}
    P_n = \{0, \epsilon/8n, \frac{1}{n} - \epsilon/16n, \frac{1}{n} + \epsilon/16n, 1\}
    \text{ for } n > 1 \text{ and } P_1 = \{0, 1\}
\end{equation*}

Then summing length times $\sup$ on each interval
\begin{align*}
    U(f_n, P_n) &= (\epsilon/8n - 0)(n) \\
    &\quad + ((\frac{1}{n} - \epsilon/16n) - (\epsilon/16n))(n) \\
    &\quad + ((\frac{1}{n} + \epsilon/16n) - (\frac{1}{n} - \epsilon/16n))(0) \\
    &= (\epsilon/8n)(n) + (\frac{1}{n} - \epsilon/16n)(n) + (\epsilon/16n)(0) \\
    &= (\epsilon/8) + (1 + \epsilon/8) \\
    &= 1 + \epsilon/4
\end{align*}
and summing lengths times $\inf$ on each interval
\begin{align*}
    L(f_n, P_n) &= (\epsilon/8n - 0)(0) \\
    &\quad + ((\frac{1}{n} - \epsilon/16n) - (\epsilon/16n))(n) \\
    &\quad + ((\frac{1}{n} + \epsilon/16n) - (\frac{1}{n} - \epsilon/16n))(n) \\
    &= (\epsilon/8n)(0) + (\frac{1}{n} - \epsilon/8n)(n) + (\epsilon/8n)(0) \\
    &= 1 - \epsilon/8
\end{align*}

Giving for every $\epsilon > 0$
\begin{align*}
    U(f_n, P_n) - L(f_n, P_n) = 3\epsilon/8
    &< \epsilon
\end{align*}

This gives
\begin{equation*}
    1 - \epsilon/8 \leq L(f_n) = \int_{0}^{1} f_n = U(f_n) \leq 1 + \epsilon/4
\end{equation*}

Since $L(f_n, P_n) \leq L(f_n)$ and $U(f_n) \leq U(f_n, P_n)$ for all paritions $P$.
Since the above equation holds for all $\epsilon > 0$ we have
\begin{equation*}
    \int_{0}^{1} f_n = 1
\end{equation*}
for all fixed n

giving
\begin{equation*}
    \lim_{n \to \infty}\int_{0}^{1} f_n = 1
\end{equation*}

Now taking any partition P will give
\begin{equation*}
    \int_{0}^{1} f = 0
\end{equation*}
Since $\sup$ and $\inf$ will be 0 on any interval as $f(x) = 0$ everywhere.

This gives
\begin{equation*}
    \lim_{n \to \infty}\int_{0}^{1} f_n = 1 \neq 0 = \int_{0}^{1} f
\end{equation*}



\end{proof}

(No matter how we integrate this its not going to work both 
riemann and lebesgue measure break here, its a motivation for DCT)



\subsection{How Lebesgue Measure Explains the failure; and then saves us through DCT}

(Show that for a function on [0, 1] enumerating the ratioanls convergeing to dirchlet
the function at the limit isnt in the class of riemann integrable functions
we know this because lebesgue defines whats riemann integrable;
then we say that under DCT which is satisifies integrating with respect to measure
we can swap the integral and the limit and it holds even though its pointwise)

\end{document}